% !TEX root = Hogwarts_1347.tex

\chapter{The Caravan}
 
A long caravan of carriages came down the narrow track from the hills, rounding the last bend in single file before drawing to a halt along the length of the pier. Thomas counted them from the rail --- more than thirty, each one dark-bodied and well-built, with large windows and iron-rimmed wheels. They had no horses. The shafts at the front of each carriage extended forward into empty air, and yet the carriages had arrived under their own momentum, their wheels turning with purpose, the spacing between them oddly generous --- far wider gaps than the arrangement seemed to require.
 
Sitting on the driver's bench of the first carriage, reins held loosely in both hands despite the apparent absence of anything to rein, was a remarkably small witch. She had short, fiery red hair --- a brighter, more untamed shade than his mother's, Thomas thought --- and even from a distance, her frame was striking in its compactness. She could not have been much more than one hundred and fifty centimetres tall. She wore a plain linen shirt tucked into trousers that had been folded twice at the hems and no shoes, her feet bare, her toes curled over the edge of the footboard as though the contact with the wood was a preference rather than an oversight.
 
The caravan drew to a halt along the length of the pier, the carriages settling with a series of soft creaks. The small witch dropped from her bench with the easy agility of a person far more comfortable outdoors than in and crossed the stone pier toward Headmaster Hepburn in a rapid, light-footed walk that was almost a run. She threw her arms around him. It was not a formal greeting; it was the embrace of a colleague who was also a dear friend, warm and unhesitating, and Hepburn returned it with a booming laugh that carried across the water.
 
\enquote{That's Professor Saoirse Millan,} said a voice at his shoulder.
 
Thomas started slightly. He hadn't heard the French sisters approach --- Celine a step behind her sister, both carrying their things.
 
\enquote{She teaches Care of Magical Creatures,} Cecile continued. \enquote{She's quite good. And a bit\ldots{} peculiar.}
 
Thomas looked back at the professor. Hepburn and Millan had begun walking together along the line of carriages, talking in low voices. The language was not Latin, and it was not English. Thomas could not place it --- it had a lilting, rounded quality, the vowels open and musical. He caught only fragments, a word here and there. One word he was almost certain he heard was \textit{selkies}. The effect on Millan was immediate. She stopped mid-stride, turned fully to face the Headmaster, and for a brief instant the easy, kinetic energy of her manner was replaced by an absolute, focused stillness --- the attention of a naturalist who has just been told about a sighting. Then she was moving again, matching Hepburn's longer stride with two of her own, firing rapid questions.
 
They reached the middle of the caravan. Hepburn drew his wand and pressed the tip firmly against his throat. When he spoke, his voice carried across the entire pier with the amplified clarity the students had grown accustomed to aboard the ships.
 
\enquote{All students are to board the carriages in groups of six. Take your cats, owl cages, and toad boxes with you inside the carriage. Leave your trunks and all other luggage on the ground in front of the carriages --- the staff will secure them on the roof. Do not attempt to load them yourselves. Don't make a mess of it.}
 
The last line drew a few laughs. Hepburn removed his wand from his throat and nodded to Millan. What followed silenced the laughter.
 
Millan drew her wand --- a short, light piece of wood with a handle of polished antler, into which the delicate, stylised forms of various creatures had been engraved with fine, loving care --- and raised it. Without a word --- not a syllable, not a murmur --- she began to move. Four carriages' worth of trunks rose simultaneously, lifting in neat, coordinated rows and drifting through the air to settle atop their respective carriages. The precision was startling --- each trunk rotated to fit flush against its neighbour, settling onto the carriage roofs with barely a sound. As the luggage found its place, leather straps zipped across the trunks lengthways, cinching tight with a series of sharp, satisfying clicks, followed immediately by a second set running perpendicular to the first, locking everything down in a firm, criss-crossed lattice. The whole operation took less than two minutes.
 
Cecile, who had spent the better part of a morning drilling \textit{Wingardium Leviosa} with the younger students on the ship and understood exactly what simultaneous levitation across four separate targets actually required, said nothing. But she exhaled slowly through her nose in a display of admiration.
 
With the last trunk secured, Millan apparated. There was no dramatic gesture, no swirl of a cloak --- one moment she was on the pier, the next she was on the driver's bench of the first carriage, settling back into her seat with her bare feet up on the footboard as though she had never left it.
 
\section*{The Land of Fire and Ice}
 
Thomas, Colum, Cecile, and Celine climbed into the fourth carriage.
 
The interior was a surprise. After the utilitarian simplicity of the \textit{Siren}'s cabins --- the narrow berths, the folding table, the heavy curtain for a door --- the carriage felt almost like entering a study. The space was compact but meticulously built. The walls were panelled in dark oak, joined with a cabinet-maker's precision, every seam flush, every corner reinforced with small brass fittings that caught the light. Two long benches faced each other across the width of the carriage, upholstered in a deep, wine-coloured wool, firm rather than soft, the kind of upholstery that would last decades without losing its shape. A narrow shelf ran along one side beneath the windows, and above it hung two small brass lamps in gimballed brackets, unlit but clean and well-maintained. The ceiling was low, vaulted slightly, and lined with the same dark oak as the walls. It was not luxurious. It was craftsmanship --- the work of a builder who believed that function and beauty were the same thing.
 
The windows were large, far larger than those on the ship, offering an unobstructed view of the world outside. Colum took the rear-facing bench on the left; Cecile sat opposite him, adjusting her robes with her usual brisk efficiency. Celine settled beside her sister, setting Pippin's cage on the floor between her feet. The cage was empty. Pippin himself was already on her lap, settled comfortably against her stomach with his eyes half-closed, leaning into the slow, methodical head-scratches she was administering with one hand. The young owl had, it appeared, traded his loyalties for a reliable supply of affection and the occasional pastry crumb, and showed no sign of renegotiating the terms. Thomas looked at the scene with an expression of theatrical offence that fooled no one, least of all himself --- the corner of his mouth kept betraying him. Thomas took the remaining rear-facing seat beside Colum and set his things down.
 
Two seats remained.
 
A few minutes later, the carriage door opened again, and two more students climbed in. First years, by the look of them, but not from the \textit{Siren} --- Thomas was certain of it. He had spent enough days in the ship's corridors and common room to recognise every face aboard, and these two were new.
 
The girl entered first. She was pale --- strikingly so, her skin almost translucent, carrying the particular quality of a person raised in a climate where the sun did not insist. Her hair was platinum blonde, cut to the shoulders, straight and fine, held back from her face by a single practical braid along one temple. Her eyes were a sharp, cold blue that took in the carriage and its occupants in a single, unhurried sweep.
 
Thomas looked at her and felt a quiet, private jolt of recognition. The paleness, the platinum hair --- the colouring was close enough to his own natural state to produce an odd, uncomfortable echo. On her, the combination was striking, it looked right. The severity of it suited her angular features and gave her an air of composed intensity that was entirely her own. It fit her beautifully. Thomas ran a hand through his dark brown hair --- an unconscious gesture --- and looked away.
 
The boy who followed her was her visual opposite in almost every regard. He was tall --- unusually so for eleven --- and powerfully built, with dark skin and a strong, defined jawline. He carried himself with a quiet confidence that was not arrogance but simple self-possession, the ease of a boy comfortable in his own frame. His eyes moved around the carriage with an appreciative attention to detail --- the brass fittings, the oak panelling, the quality of the joinery --- that suggested a person who noticed how things were made.
 
They spoke to each other as they settled in with the relaxed familiarity of people who had shared close quarters for an extended period --- many days inside a small ship, most likely. Their Latin was fluent but distinctly accented --- his carrying a warm, rounded quality, hers clipped and precise, the vowels shorter, the consonants harder. Neither accent was British.
 
The boy spoke first, addressing the carriage with a natural warmth.
 
\enquote{Efe Adeyemi,} he said, extending a broad hand toward Thomas, the nearest to him. His grip was firm and brief. He worked around the carriage --- Colum, Cecile, Celine --- greeting each with an easy directness.
 
\enquote{Astrid Ragnarsdóttir,} the girl said. Her voice was lower than Thomas had expected, steady and measured. She nodded to the group rather than offering a hand, but the gesture carried its own courtesy.
 
Celine, whose curiosity operated on a schedule that did not wait for appropriate moments, tilted her head at the name. \enquote{Ragnarsdóttir,} she repeated carefully, her French accent wrapping around the unfamiliar syllables. \enquote{Where is your family name from?}
 
\enquote{It isn't a family name,} Astrid said. She did not correct Celine with any sharpness --- it was simply a fact that needed stating. \enquote{It's a patronymic. It means daughter of Ragnar. My father is Ragnar, so I am Ragnarsdóttir. It's the custom in Iceland.}
 
\enquote{The land of fire and ice,} completed Efe.
 
Celine absorbed this with visible fascination, the kind that would almost certainly produce follow-up questions at a later and less convenient time.
 
She turned to Efe. \enquote{And where are you from?}
 
\enquote{Southwark,} Efe said.
 
There was a beat. Then Efe grinned --- a warm, broad expression that confirmed the joke before anyone had to ask. Colum exhaled a short laugh. Efe's grin widened before he settled into a more measured tone.
 
\enquote{My family moved to London when I was six. From the Kingdom of Edo, in the west of Africa. We were brought to England by a merchant of rare magical items, Mr. Fox.} He said the name with a quiet, settled warmth, the way you speak of a person who had earned a permanent place in your family's story. \enquote{He offered my father a position, and my father accepted.}
 
Thomas said nothing, but the name landed with a sharp, private clarity. Mr. Fox. Elias Robbert Fox --- the crooked man in the room behind the red door, leaning on his table, his body scarred and his smile pulling against the damage. The house-elves in their white mantles. The leather purse of coins. And standing guard at the entrance, the dark-skinned wizard with the magnificent African blackwood wand carved with minute figurines and inlaid with bronze --- Efe's father, Thomas now realised. He filed the connection away without comment.
 
\section*{The Glen}
 
The carriage lurched gently and began to move.
 
For the first stretch of the journey, the conversation continued --- the easy, exploratory talk of people recently introduced, establishing the basic coordinates of who they were and where they had come from. Efe and Astrid, it emerged, had traveled together aboard the \textit{Melusine}, the rearmost ship of the convoy, and their rapport had the quality of a friendship forged in confined quarters over shared meals and long days at sea.
 
But gradually, as the road climbed and the landscape opened around them, the talk thinned and stopped. Not because anything had gone wrong between them, but because the windows demanded it.
 
The country on either side of the narrow track was unlike anything Celine had ever seen. The road ran high along the shoulder of a glen, following the curve of a long, dark loch that lay far below them, its surface black and perfectly still, reflecting the mountains on its far shore with such fidelity that the reflection was indistinguishable from the thing itself. The hills rose steeply from the water's edge --- not the gentle, rolling green of southern England but harder, their flanks covered in a dense, dark pelt of oak and birch and Scots pine that gave way, higher up, to bare rock and heather the colour of rusted iron.
 
Where the glen narrowed, the mountains pressed close, their ridges running parallel to the road like the walls of a vast, roofless corridor. Waterfalls cut white threads down the dark rock faces, some dropping from heights so great that the water dissolved into mist before it reached the bottom. In the hollows between the peaks, pockets of low cloud sat motionless, as though the mountains had caught them and were keeping them.
 
The light was different here. Far above the trapped mist in the hollows, the higher cloud cover was in constant motion, and the sun found its way through the gaps in long, horizontal shafts, illuminating single stretches of hillside in vivid, burning gold while the land on either side remained in deep shadow. The effect shifted constantly --- a ridge would blaze for a moment and then darken, and the light would find a new slope, a new angle, a new arrangement of colour that lasted only seconds before changing again.
 
Thomas leaned against the window frame and watched it pass. He had no memory of this place, not exactly. He had been too young when they left for London to carry detailed pictures in his mind --- no specific ridge, no particular tree, no view that he could point to and say \textit{I remember that}. But the impression was there, deep and inarticulate: the quality of the air, the scale of the silence, the way the land felt --- not welcoming, not hostile, but present, aware of you in the way that very old places are. It settled over him like a scent he could not name but knew he had encountered before.
 
He wished the caravan moved more slowly. The beauty outside the windows deserved more time than the road was giving it.
 
Then Cecile, who had been sitting facing forward, straightened in her seat.
 
\enquote{There,} she said quietly.
 
Thomas turned.
 
Very far ahead, still small with distance but already unmistakable against the sky, the castle sat upon a high, dark outcrop of rock above the shore of a vast lake. It rose from the stone as though it had been grown from it rather than built upon it, a sprawling mass of towers and turrets and high curtain walls, the stonework catching the afternoon light along its western face. The towers were too many to count at this distance --- some slender and impossibly tall, others broad and fortified, their crenellated battlements sharp against the grey sky. Windows glinted across the façade in scattered points of reflected light, hundreds of them, as though the castle were watching the valley below with a thousand bright, quiet eyes.
 
Thomas, Colum and Astrid rose from their seats and moved to the front-facing bench, pressing in beside  Cecile, Celine and Efe so that all six of them could watch the castle through the back window as the caravan continued its approach. The bench was not built for six, and the arrangement required a degree of compression that would have been uncomfortable if anyone had been thinking about comfort. No one was.
 
Thomas looked at Celine. She had gone entirely still. Her lips were slightly parted, her eyes wide, the half-eaten apple in her hand forgotten. She was not trying to analyse what she was seeing or compare it to anything she had known. She was simply receiving it --- the whole vast, impossible, luminous fact of the place --- with the unguarded awe of a girl who had traveled a very long way from a bakery in Lyon and was now looking at the castle where her life was about to change. She did not say a word. She didn't need to.
 
\section*{The Painted Posts}
 
As the road descended from the high glen and the distance to the castle closed, the landscape on both sides of the narrow track shifted. The wild, uninhabited country of the mountains gave way --- gradually and then all at once --- to the marks of habitation.
 
On the left, set back from the road at the bottom of a gentle slope, the village of Hogsmeade --- the only all-wizarding village in the British Isles --- came into view.
 
It was not large, but it was deeply, unmistakably itself. The buildings were a crooked, comfortable collection of stone and timber, their roofs thatched or slated in dark grey, their chimneys --- many of them slightly off-vertical, as though the houses had settled into the field over the years and brought their chimneys with them --- trailing thin threads of smoke in colours that no Muggle fire would produce. The lanes between the buildings were narrow and winding, following no grid, respecting no surveyor's line, running instead along paths that had been worn into the ground by feet and hooves and the habits of many generations. Shop fronts faced the main street, their signs hand-painted and slightly crooked, and through the carriage window Thomas caught glimpses of warm light behind thick glass, the shapes of people moving inside.
 
Thomas looked at it. The view unlocked nothing specific --- no street, no doorway, no face. But it unlocked a feeling: a low, formless warmth that sat behind his ribs, the residue of a place that had once been home to a very small boy and a family that was still whole. He did not try to sharpen it into a memory. He let it sit where it was.
 
On the right, the mountains fell away entirely, opening onto a broad, flat field that stretched toward a dark lake at its far edge. The field was vast --- far larger than any sporting ground Thomas had seen in London --- and it was occupied.
 
At the near end stood a large oval arena built entirely of heavy timber, its structure open to the sky. Two grandstands ran along the longer sides, their seats rising in steep tiers. Beyond the arena, scattered across the field at irregular intervals, stood smaller installations --- targeting ranges, jousting tracks, and obstacle courses of various configurations. But what dominated the field entirely were the posts.
 
A series of massive wooden columns --- enormously tall, painted in vivid, saturated colours that stood out against the green and grey of the landscape --- rose from the ground at varying heights and distances from one another. They were everywhere. They ran the full length of the field, from the shallows of the lake at the far edge, where their reflections doubled their height in the still water, through the open ground past the arena, and all the way to the near side, close enough to the road that Thomas could see the grain of the painted timber. The winding course they described snaked between the other installations, around natural features of the terrain, and back again in long, sweeping loops. It was, unmistakably, a racing circuit.
 
\enquote{That's broom racing,} Cecile said, leaning forward to look past Celine. \enquote{Course racing, specifically. And that ---}
 
A figure had launched from the far end of the circuit. It moved with a speed of an arrow loosed from a longbow, a sheer velocity that seemed to ignore the air rather than negotiate with it. The figure reached the first post, banked hard, and came around it in a turn so tight that the angle defied sense.
 
 
\enquote{--- must be professor da Senna,} Cecile finished.
 
The rider flew closer, contouring a post near the road, and Thomas watched carefully. He did not slow as the turn approached. If anything, the broom accelerated on the exit of the turn, the rider leaning into the curve and then surging out of it with a burst of speed that left a faint, visible disturbance in the air behind it. Then he was gone, banking behind a distant post, and the circuit was empty again, and the field was quiet.
 
\section*{The Winged Boars}
 
The caravan slowed.
 
Ahead, the road narrowed to a single track and approached a tall iron gate set between stone pillars. Atop each pillar sat a carved stone boar, winged and rearing, worn smooth by weather but still proud. The gate itself was elegant --- wrought iron, ornamented rather than fortified. It looked, Thomas thought, considerably more fragile than the entrance to a castle of this size ought to be. But then, the real protections of a place like this would not be the kind you could see.
 
Professor Millan dropped from the driver's bench of the first carriage and walked toward the gate. She drew her wand, raised it, and began. The first movements were small --- quick, precise flicks directed at specific points along the gate's framework, as though she were unlocking a sequence of invisible catches. Then the gestures widened, the wand tracing broader arcs through the air, her arm extending fully. She did not speak. Not a word, not a whisper. The spells were entirely silent, and yet their effect was immediate: a translucent barrier shimmered into visibility across the full span of the gateway, hanging in the air like a sheet of still water suspended on its edge. It held for a moment --- long enough for Thomas to see the faint distortions it produced in the landscape behind it, the trees and the sky bending slightly as though seen through old glass --- and then it parted, drew aside, and folded into nothing.
 
With a final, slow flick, the iron gate began to swing open. It moved with a ponderous weight that contradicted its appearance entirely, as though the metal were ten times heavier than it looked, the hinges resisting and then yielding with a deep, resonant groan that Thomas felt in his chest.
 
The first carriage rolled through. Then the second. One by one, the carriages passed between the stone pillars and their winged boars, moving slowly, the track beyond the gate curving gently upward through a corridor of tall trees.
 
When the last carriage had passed through, Millan closed the gate. It swung shut with the same heavy, slow authority with which it had opened. She raised her wand again, and the sequence reversed --- the translucent barrier reappeared, shimmering into existence across the gateway, holding for a moment, and then fading until the gate looked exactly as it had before: elegant, ornamental, and entirely insufficient for the task of protecting what lay beyond it.
 
Millan tucked her wand into her belt and walked back along the line of carriages toward the front, her bare feet silent on the packed earth of the road. She climbed back onto her bench, settled into her seat, and the caravan moved forward through the final stretch of tree-lined road toward the castle.
 
The towers filled the sky above the canopy now, close enough to see the individual stones of the walls and the shapes of the gargoyles perched along the parapets. The sound of the carriage wheels changed as the track gave way to paved stone. The trees thinned, and then opened, and then fell away entirely, and Hogwarts castle rose before them in its full, staggering immensity --- closer, more real and more overwhelming than it had been from any distance.
 
The caravan slowed to a walk, and then stopped.
 
